\section*{Chapitre \ref{ch:m3:commodity-options-and-structured-hedges} --- Options sur matières premières et couvertures structurées}

\begin{solution}{exo:m3:commodity-options-and-structured-hedges:1}
Le swap verse $62 - 55 = \$7$ par baril~; le producteur réalise 62~\$.
\end{solution}

\begin{solution}{exo:m3:commodity-options-and-structured-hedges:2}
La moyenne des prix sur la période est moins dispersée que le prix à son terme, car les fixations précoces
ne peuvent plus bouger une fois fixées, et la moyenne lisse le reste~; un sous-jacent moins volatil rend
l'option moins chère.
\end{solution}

\begin{solution}{exo:m3:commodity-options-and-structured-hedges:3}
50~\$, 60~\$ et 75~\$.
\end{solution}

\begin{solution}{exo:m3:commodity-options-and-structured-hedges:4}
45~\$ (le prix plus le put spread de 15 dollars), 60~\$, 65~\$ et 70.24~\$.
\end{solution}

\begin{solution}{exo:m3:commodity-options-and-structured-hedges:5}
$35\%/\sqrt3 = 20.2\%$ vaut pour une moyenne continue~; avec douze fixations, la dernière à l'échéance, le
facteur exact est $\sqrt{(n+1)(2n+1)/(6n^2)} = 0.613$, soit 21.5~\%, et la valeur de la simulation, 21.3~\%, en est proche (le
logarithme d'une moyenne n'est pas exactement normal).
\end{solution}

\begin{solution}{exo:m3:commodity-options-and-structured-hedges:6}
Il peut acheter la matière première pour la date proche et la vendre pour la date éloignée chaque fois que le spread dépasse son
coût de stockage, et ne rien faire sinon~: son gain est un call sur le spread entre date éloignée et date proche,
de strike égal au coût de stockage.
\end{solution}

\begin{solution}{exo:m3:commodity-options-and-structured-hedges:7}
Call moins put vaut $10.31 - 5.51 = 4.80 = e^{-0.04}(60 - 55)$~; le call vaut 10.31~\$.
\end{solution}

\begin{solution}{exo:m3:commodity-options-and-structured-hedges:8}
Elle ne protège qu'entre 60~\$ et le strike du put le plus bas~: en dessous, le producteur suit le prix à la baisse avec un
coussin fixe, justement dans l'effondrement qu'il voulait assurer. Elle coûte aussi son gain au-dessus du strike du call.
«~Ne coûte rien~» signifie seulement que les primes se compensent aujourd'hui.
\end{solution}

\begin{solution}{pb:m3:commodity-options-and-structured-hedges:1}
\textbf{1.} 2.66~\$ par baril. \textbf{2.} Environ 664~millions de dollars. \textbf{3.} 4.4~\% de 15~milliards de dollars.
\textbf{4.} 2.5~milliards de dollars. \textbf{5.} 6.25~milliards de dollars. \textbf{6.} 70.24~\$. \textbf{7.} 3.75~milliards de dollars.
\textbf{8.} Même montant, 3.75~milliards de dollars~: la protection s'arrête à 45~\$. \textbf{9.} Environ 4.94~milliards de dollars versés au
call. \textbf{10.} Le \omterm{def:m3:commodity-options-and-structured-hedges:collar}{three-way collar} fait mieux lors des baisses modérées (jusqu'à environ 45~\$), le put lors des effondrements et des
hausses. \textbf{11.} Le budget dépend du prix moyen d'exportation de l'année, et les \omterm{def:m3:commodity-options-and-structured-hedges:apo}{options sur prix moyen} sont moins chères.
\textbf{12.} Ils vendent des contrats à terme (et ajustent au fil des fixations de la moyenne), ce qui ajoute une pression vendeuse lorsque le
programme est placé et à mesure que les prix baissent. \textbf{13.} Pour limiter l'impact sur le marché d'un acheteur connu et de grande taille, ainsi que des
couvertures des banques. \textbf{14.} L'excédent de protection devient une position spéculative~; le ministère peut vendre
les puts excédentaires ou accepter le versement. \textbf{15.} Comme coût d'assurance, réparti sur l'année qu'il couvre.
\textbf{16.} Les puts plafonnent le risque de baisse du budget tout en conservant la hausse, et coûtent un montant connu~; un swap abandonne
la hausse et peut exiger des versements en cas de hausse. \textbf{17.} De gros versements les années de crise et des primes perdues
les autres années~: sa valeur réside dans la certitude budgétaire qu'elle procure, et non dans un profit. \textbf{18.} Lorsque le risque budgétaire tient à un
effondrement profond, que le put vendu laisse non couvert, ou lorsque le gain en cas de hausse compte. \textbf{19.}
\emph{Résultat nommé~:} le put coûte 4.4~\% du revenu budgété (664~millions de dollars) et verse 2.5~milliards de dollars pour une
année à 45~\$ et 6.25~milliards de dollars pour une année à 30~\$, contre 3.75~milliards de dollars dans les deux cas pour le \omterm{def:m3:commodity-options-and-structured-hedges:collar}{three-way collar} de même
prime. \textbf{20.} Un floor sous son budget, à prime connue, tout en conservant la hausse.
\end{solution}

\begin{solution}{iq:m3:commodity-options-and-structured-hedges:1}
Une option sur la moyenne d'un prix sur une période~: des producteurs et des consommateurs qui achètent ou vendent chaque jour et
sont exposés à la moyenne, et non au prix d'une seule journée.
\iqlookfor{adéquation à l'exposition, et faible coût.}
\end{solution}

\begin{solution}{iq:m3:commodity-options-and-structured-hedges:2}
Vous achetez un floor et le financez en vendant un cap~: votre prix restera compris entre les deux, quel que soit le
comportement du marché, et rien n'est payé d'avance. Vous renoncez aux gains au-delà du cap.
\iqlookfor{langage clair et simple sur le floor, le cap et le coût d'opportunité.}
\end{solution}

\begin{solution}{iq:m3:commodity-options-and-structured-hedges:3}
La moyenne d'une trajectoire brownienne sur $[0,T]$ a pour variance $\int_0^T\!\int_0^T \min(s,u)\,ds\,du/T^2 = T/3$,
donc son écart-type vaut $\sigma\sqrt{T/3}$~: $1/\sqrt3$ de celui de la valeur terminale.
\iqlookfor{l'intégrale de la covariance du mouvement brownien.}
\end{solution}

\begin{solution}{iq:m3:commodity-options-and-structured-hedges:4}
Vente d'un put~: on achète des contrats à terme pour être delta-neutre (on en vend si la position du client rend long en delta),
on réduit la couverture au fur et à mesure des fixations (les fixations passées n'ont pas de delta), on gère le gamma et le vega avec des
options cotées, et on anticipe l'impact de marché de la couverture.
\iqlookfor{delta qui décroît avec les fixations, gestion du vega, liquidité.}
\end{solution}

\begin{solution}{iq:m3:commodity-options-and-structured-hedges:5}
Un contrat qui permet au détenteur de faire varier ses prélèvements quotidiens entre des bornes, soumis à des plafonds annuels et à une
\omterm{def:m3:commodity-options-and-structured-hedges:swing}{clause take-or-pay} minimale. La flexibilité tient à un ensemble d'options permettant de prélever davantage quand le prix est élevé et moins quand il est
bas, ce que ne possède pas un contrat à volume fixe.
\iqlookfor{flexibilité vue comme des options, et contraintes annuelles.}
\end{solution}

\begin{solution}{iq:m3:commodity-options-and-structured-hedges:6}
Données d'entrée~: prévisions de production assorties d'incertitude, seuils budgétaires et de covenant, courbes de marché et smiles,
lignes de crédit. Structures~: swaps, puts, \omterm{def:m3:commodity-options-and-structured-hedges:collar}{collars}, \omterm{def:m3:commodity-options-and-structured-hedges:collar}{three-way collars} sur des périodes choisies. Sorties~: coût, graphiques de
gains, distribution des flux de trésorerie et probabilité de non-respect du covenant selon des prix et volumes simulés, et
besoins de marge.
\iqlookfor{objectif lié aux covenants, risque conjoint de prix et de volume.}
\end{solution}
